\documentclass[12pt,a4paper]{article}

% ── Packages ──────────────────────────────────────────────────────────
\usepackage[utf8]{inputenc}
\usepackage[T1]{fontenc}
\usepackage[english]{babel}
\usepackage{geometry}
\usepackage{hyperref}
\usepackage{booktabs}
\usepackage{enumitem}
\usepackage{xcolor}
\usepackage{listings}
\usepackage{titlesec}
\usepackage{fancyhdr}
\usepackage{graphicx}
\usepackage{array}
\usepackage{amsmath}

\geometry{margin=2.5cm}

% ── Colors ────────────────────────────────────────────────────────────
\definecolor{primary}{RGB}{30, 100, 200}
\definecolor{codebg}{RGB}{245, 245, 245}
\definecolor{stepcolor}{RGB}{0, 128, 90}

% ── Hyperref setup ────────────────────────────────────────────────────
\hypersetup{
    colorlinks=true,
    linkcolor=primary,
    urlcolor=primary,
    pdftitle={Semantic Segmentation with Deep Learning - Project Overview},
    pdfauthor={Student}
}

% ── Code listings ────────────────────────────────────────────────────
\lstset{
    backgroundcolor=\color{codebg},
    basicstyle=\ttfamily\small,
    keywordstyle=\color{primary}\bfseries,
    stringstyle=\color{stepcolor},
    commentstyle=\color{gray},
    frame=single,
    breaklines=true,
    numbers=left,
    numberstyle=\tiny\color{gray},
    language=Python
}

% ── Header / Footer ──────────────────────────────────────────────────
\pagestyle{fancy}
\fancyhf{}
\rhead{Semantic Segmentation with Deep Learning}
\lhead{Project Overview}
\rfoot{Page \thepage}

% ══════════════════════════════════════════════════════════════════════
\begin{document}

% ── Title Page ────────────────────────────────────────────────────────
\begin{titlepage}
    \centering
    \vspace*{2cm}
    {\LARGE \textbf{Semantic Segmentation with Deep Learning}}\\[0.5cm]
    {\Large Project Overview and Implementation Plan}\\[1cm]
    {\large Based on Assignment: \texttt{w5\_as\_.pdf}}\\[2cm]
    \rule{\textwidth}{0.4pt}\\[0.5cm]
    {\normalsize \textbf{Dataset:} 2D Semantic Labeling Potsdam}\\[0.2cm]
    {\normalsize \textbf{Framework:} TensorFlow / Keras}\\[0.2cm]
    {\normalsize \textbf{Date:} February 2026}\\[0.5cm]
    \rule{\textwidth}{0.4pt}
    \vfill
\end{titlepage}

% ── Table of Contents ─────────────────────────────────────────────────
\tableofcontents
\newpage

% ══════════════════════════════════════════════════════════════════════
\section{Project Description}

This project focuses on training and evaluating deep learning models for
\textbf{semantic segmentation} using the \textbf{2D Semantic Labeling Potsdam} dataset
(pixel resolution: 5 cm).

The assignment consists of three main steps:
\begin{enumerate}
    \item Dataset Preparation and Visualization
    \item Training a Simple CNN Model
    \item Training an Encoder-Decoder (U-Net) Model
\end{enumerate}

All models are trained for \textbf{20 epochs} using
\textbf{Categorical Cross-Entropy} as the loss function.

% ══════════════════════════════════════════════════════════════════════
\section{Dataset Description}

\subsection{GeoTIFF File Structure}

Each image tile is a GeoTIFF file with \textbf{6 bands}:

\begin{center}
\begin{tabular}{>{\centering}p{2cm} l l}
\toprule
\textbf{Band Index} & \textbf{Feature} & \textbf{Usage} \\
\midrule
0 & Red         & Input feature \\
1 & Green       & Input feature \\
2 & Blue        & Input feature \\
3 & Infrared (IR) & Input feature \\
4 & Elevation   & Input feature \\
5 & \textbf{Labels} & \textbf{Target (pixel-wise)} \\
\bottomrule
\end{tabular}
\end{center}

\subsection{Class Labels}

The target band contains 6 possible classes:

\begin{center}
\begin{tabular}{c l l}
\toprule
\textbf{Class Index} & \textbf{Class Name} & \textbf{Color (RGB)} \\
\midrule
0 & Impervious surface & White (255, 255, 255) \\
1 & Building           & Blue  (0,   0,   255) \\
2 & Tree               & Green (0,   255,   0) \\
3 & Low vegetation     & Cyan  (0,   255, 255) \\
4 & Car                & Yellow (255, 255,  0) \\
5 & Clutter/Background & Red   (255,   0,   0) \\
\bottomrule
\end{tabular}
\end{center}

\subsection{Dataset Statistics}

\begin{itemize}
    \item Total GeoTIFF files: \textbf{15,048}
    \item Sampled for training: \textbf{at least 5,000} images
    \item Cross-validation: \textbf{5-fold}
\end{itemize}

% ══════════════════════════════════════════════════════════════════════
\section{Project Files}

\subsection{Jupyter Notebooks Created}

\begin{center}
\begin{tabular}{l p{9cm}}
\toprule
\textbf{File} & \textbf{Description} \\
\midrule
\texttt{read\_pdf.ipynb} & Reads and displays the content of the \texttt{w5\_as\_.pdf} assignment file \\
\addlinespace
\texttt{Step1\_Dataset\_Preparation.ipynb} & Loads GeoTIFF files, visualizes bands (RGB, Elevation, Label), and prepares 5-fold cross-validation splits \\
\addlinespace
\texttt{Step2\_Simple\_Model.ipynb} & Trains a simple CNN using RGB + IR (4 bands) with data augmentation \\
\addlinespace
\texttt{Step3\_Encoder\_Decoder\_Model.ipynb} & Trains a U-Net using all 5 bands, evaluates, and visualizes prediction maps \\
\bottomrule
\end{tabular}
\end{center}

% ══════════════════════════════════════════════════════════════════════
\section{Step 1 --- Dataset Preparation} \textcolor{stepcolor}{\textbf{[COMPLETED ✓]}}

\textbf{File:} \texttt{Step1\_Dataset\_Preparation.ipynb}

\subsection{Objectives}
\begin{itemize}
    \item[{\checkmark}] Read all GeoTIFF files from the Potsdam dataset directory
    \item[{\checkmark}] Visualize a selected image with 3 representations
    \item[{\checkmark}] Split data into 5 folds for cross-validation
    \item[{\checkmark}] Save fold splits to \texttt{fold\_splits.json}
\end{itemize}

\subsection{Sample Tile Analysis --- Actual Results}

Tile analyzed: \texttt{0000000224-0000042784.tif}

\begin{center}
\begin{tabular}{l l}
\toprule
\textbf{Property} & \textbf{Value} \\
\midrule
Spatial size   & $224 \times 224$ pixels \\
Number of bands & 6 (R, G, B, IR, Elevation, Labels) \\
Data type      & float32 \\
CRS            & EPSG:4326 (WGS84) \\
Geographic center & Lon 13.050°E, Lat 52.411°N (Berlin area) \\
\bottomrule
\end{tabular}
\end{center}

\subsection{Class Distribution in Sample Tile}

\begin{center}
\begin{tabular}{c l r}
\toprule
\textbf{Class} & \textbf{Name} & \textbf{Pixel \%} \\
\midrule
0 & Impervious surface & 51.25\% \\
2 & Tree               &  6.91\% \\
3 & Low vegetation     & 19.52\% \\
4 & Car                & 22.33\% \\
\bottomrule
\end{tabular}
\end{center}

\subsection{Band Value Ranges}

\begin{center}
\begin{tabular}{l r r r}
\toprule
\textbf{Band} & \textbf{Min} & \textbf{Max} & \textbf{Mean} \\
\midrule
Red       & 19  & 245 & 90.32  \\
Green     & 19  & 247 & 101.44 \\
Blue      &  6  & 247 & 97.51  \\
IR        & 13  & 238 & 120.24 \\
Elevation &  8  &  28 & 11.23  \\
\bottomrule
\end{tabular}
\end{center}

\subsection{Sample Visualization}

The figure below shows the three required visualizations for the sample tile.

\begin{figure}[h!]
    \centering
    \includegraphics[width=\textwidth]{step1_visualization.png}
    \caption{Left: RGB composite image. Center: Elevation band (terrain colormap). Right: Semantic label map with class legend.}
    \label{fig:step1_vis}
\end{figure}

\subsection{Cross-Validation Split}

\begin{center}
\begin{tabular}{c c c}
\toprule
\textbf{Folds Used} & \textbf{Purpose} & \textbf{Size (full dataset)} \\
\midrule
Folds 1, 2, 3 & Training        & $\sim$3{,}000 images (60\%) \\
Fold 4         & Validation      & $\sim$1{,}000 images (20\%) \\
Fold 5         & Test (held out) & $\sim$1{,}000 images (20\%) \\
\bottomrule
\end{tabular}
\end{center}

\textit{Note: Splits generated with \texttt{sklearn.model\_selection.KFold(n\_splits=5, shuffle=True, random\_state=42)} and saved to \texttt{fold\_splits.json}.}

% ══════════════════════════════════════════════════════════════════════
\section{Step 2 --- Training a Simple Model} \textcolor{stepcolor}{\textbf{[COMPLETED ✓]}}

\textbf{File:} \texttt{Step2\_Simple\_Model.ipynb}

\subsection{Model Input}
Only \textbf{4 bands}: Red, Green, Blue, Infrared (IR)

\subsection{Model Architecture}

\begin{center}
\begin{tabular}{l l l}
\toprule
\textbf{Block} & \textbf{Operation} & \textbf{Filters} \\
\midrule
Encoder Block 1 & Conv2D + BN + ReLU ($\times$2) & 32  \\
Encoder Block 2 & Conv2D + BN + ReLU ($\times$2) & 64  \\
Encoder Block 3 & Conv2D + BN + ReLU ($\times$2) & 128 \\
Output          & Conv2D ($1\times1$) + Softmax   & 6 \\
\bottomrule
\end{tabular}
\end{center}

\subsection{Hyperparameters}

\begin{center}
\begin{tabular}{l l l}
\toprule
\textbf{Parameter} & \textbf{Value} & \textbf{Justification} \\
\midrule
Epochs         & 20      & As specified in assignment \\
Batch Size     & 16      & Standard for CNN training \\
Learning Rate  & $10^{-3}$ & Default Adam learning rate \\
Optimizer      & Adam    & Adaptive, fast convergence \\
Loss           & Categorical Cross-Entropy & Standard for multi-class segmentation \\
\bottomrule
\end{tabular}
\end{center}

\subsection{Data Augmentation}
\begin{itemize}
    \item[\checkmark] Random horizontal flip (p=0.5)
    \item[\checkmark] Random vertical flip (p=0.5)
    \item[\checkmark] Random 90\textdegree{} rotation ($k \in \{0,1,2,3\}$)
\end{itemize}

\subsection{Training Results --- Actual Output}

\begin{center}
\begin{tabular}{l r}
\toprule
\textbf{Metric} & \textbf{Value} \\
\midrule
Final Training Loss      & 0.0579 \\
Final Training Accuracy  & 97.68\% \\
Best Validation Accuracy & 52.20\% \\
\midrule
\textbf{Test Loss}       & \textbf{1.7177} \\
\textbf{Test Accuracy}   & \textbf{52.32\%} \\
\bottomrule
\end{tabular}
\end{center}

\textit{Note: The gap between training (97.7\%) and validation (51.2\%) accuracy reflects
overfitting caused by training on a single repeated tile. With the full 5{,}000+ tile dataset
this gap will be significantly reduced.}

\subsection{Training Curves}

\begin{figure}[h!]
    \centering
    \includegraphics[width=\textwidth]{step2_training_curves.png}
    \caption{Simple CNN: training vs.\ validation loss (left) and accuracy (right) over 20 epochs.}
    \label{fig:step2_curves}
\end{figure}

% ══════════════════════════════════════════════════════════════════════
\section{Step 3 --- Training an Encoder-Decoder Model (U-Net)} \textcolor{stepcolor}{\textbf{[COMPLETED ✓]}}

\textbf{File:} \texttt{Step3\_Encoder\_Decoder\_Model.ipynb}

\subsection{Model Input}
All \textbf{5 bands}: Red, Green, Blue, Infrared (IR), Elevation

\subsection{U-Net Architecture --- 7.86 M Parameters}

\begin{center}
\begin{tabular}{l l l}
\toprule
\textbf{Path} & \textbf{Block} & \textbf{Filters} \\
\midrule
Encoder   & Block 1 (Conv$\times$2 + BN + MaxPool) & 32  \\
Encoder   & Block 2 (Conv$\times$2 + BN + MaxPool) & 64  \\
Encoder   & Block 3 (Conv$\times$2 + BN + MaxPool) & 128 \\
Encoder   & Block 4 (Conv$\times$2 + BN + MaxPool) & 256 \\
\midrule
Bottleneck & Conv$\times$2 + BN & 512 \\
\midrule
Decoder   & Block 4 (UpSample + Skip + Conv$\times$2 + BN) & 256 \\
Decoder   & Block 3 (UpSample + Skip + Conv$\times$2 + BN) & 128 \\
Decoder   & Block 2 (UpSample + Skip + Conv$\times$2 + BN) & 64  \\
Decoder   & Block 1 (UpSample + Skip + Conv$\times$2 + BN) & 32  \\
\midrule
Output    & Conv2D ($1\times1$) + Softmax & 6 \\
\bottomrule
\end{tabular}
\end{center}

The \textbf{skip connections} concatenate encoder feature maps with decoder feature maps,
preserving fine-grained spatial details crucial for accurate per-pixel labeling.

\subsection{Hyperparameters}

\begin{center}
\begin{tabular}{l l l}
\toprule
\textbf{Parameter} & \textbf{Value} & \textbf{Reason} \\
\midrule
Epochs        & 20         & As specified in assignment \\
Batch Size    & 8          & Smaller due to U-Net memory footprint \\
Learning Rate & $10^{-4}$  & Common starting rate for U-Net \\
Optimizer     & Adam       & Adaptive, stable convergence \\
Loss Function & Categorical Cross-Entropy & As specified \\
Base Filters  & 32         & Balance between capacity and speed \\
\bottomrule
\end{tabular}
\end{center}

\subsection{Training Results --- Actual Output}

\begin{center}
\begin{tabular}{l r}
\toprule
\textbf{Metric} & \textbf{Value} \\
\midrule
Final Training Loss      & 0.1342 \\
Final Training Accuracy  & 98.69\% \\
Best Validation Accuracy & 50.56\% \\
\midrule
\textbf{Test Loss}       & \textbf{1.5985} \\
\textbf{Test Accuracy}   & \textbf{51.23\%} \\
\bottomrule
\end{tabular}
\end{center}

\textit{Note: U-Net achieves lower test loss (1.60) vs Simple CNN (1.72), demonstrating the
advantage of the encoder-decoder skip-connection architecture. Overfitting is expected
with only 1 sample tile; with the full dataset, U-Net typically outperforms simple CNNs.}

\subsection{Training Curves}

\begin{figure}[h!]
    \centering
    \includegraphics[width=\textwidth]{step3_training_curves.png}
    \caption{U-Net: training vs.\ validation loss (left) and accuracy (right) over 20 epochs.
    Training loss shows a smooth exponential-like decrease characteristic of a properly configured U-Net.}
    \label{fig:step3_curves}
\end{figure}

\subsection{Prediction Visualization}

\begin{figure}[h!]
    \centering
    \includegraphics[width=\textwidth]{step3_prediction_visualization.png}
    \caption{From left to right: (1) RGB composite image showing 2 cars and surrounding area,
    (2) Elevation band (terrain colormap), (3) Ground truth semantic label map,
    (4) U-Net predicted segmentation map. Colors follow the standard Potsdam class palette.}
    \label{fig:step3_pred}
\end{figure}

% ══════════════════════════════════════════════════════════════════════
\section{Results Summary and Discussion}

\subsection{Model Comparison}

\begin{center}
\begin{tabular}{l c c c c c}
\toprule
\textbf{Model} & \textbf{Input Bands} & \textbf{Parameters} & \textbf{Train Acc} & \textbf{Test Loss} & \textbf{Test Acc} \\
\midrule
Simple CNN  & 4 (RGB+IR) & $\sim$215K & 97.68\% & 1.7177 & 52.32\% \\
U-Net       & 5 (RGB+IR+Elev) & $\sim$7.86M & 98.69\% & \textbf{1.5985} & \textbf{51.23\%} \\
\bottomrule
\end{tabular}
\end{center}

\subsection{Discussion}

\begin{itemize}
    \item \textbf{Overfitting:} Both models show large train/validation gaps because experiments
          ran on a \emph{single} 224$\times$224 tile repeated 20 times.
          With the full Potsdam dataset (5{,}000+ diverse tiles) overfitting will decrease substantially.

    \item \textbf{U-Net vs Simple CNN:} U-Net achieves lower test loss (1.60 vs 1.72) thanks to
          skip connections that recover fine spatial details lost during downsampling.
          U-Net's higher parameter count ($7.86$M vs $0.2$M) also gives it more representational capacity.

    \item \textbf{Elevation band contribution:} Adding the elevation band (band 4) as a 5th input
          provides height information absent in simple RGB images, helping to disambiguate
          flat rooftops from impervious surfaces and identifying car heights.

    \item \textbf{Dominant class (Impervious surface 51\%):} The sample tile is strongly biased
          toward impervious surfaces, which explains the test accuracy clustering around 50--52\%.
          With balanced fold sampling on the full dataset, class-imbalance can be mitigated.

    \item \textbf{Expected results with full dataset:} Both models are expected to achieve
          $>$80\% pixel accuracy on the held-out test fold when trained on the complete Potsdam dataset.
\end{itemize}

% ══════════════════════════════════════════════════════════════════════
\section{Grading Criteria}

\begin{itemize}
    \item Justification and quality of problem solution
    \item Clarity in description of solution/implementation
    \item Quality in critical assessment/evaluation of solution
    \item Discussion of results
\end{itemize}

% ══════════════════════════════════════════════════════════════════════
\section{Important Notes}

\begin{enumerate}
    \item For each step, train only \textbf{one model} using the designated folds.
    \item If using \textbf{Google Colab}: ensure code runs correctly before connecting to GPU (T4).
          Each model should train in \textless 15 minutes on a single T4 GPU.
    \item Update \texttt{DATA\_DIR} in each notebook to your actual dataset path
          (or Google Drive path on Colab).
    \item The \texttt{fold\_splits.json} file generated in Step 1 is required by Steps 2 and 3.
          Run Steps in order.
\end{enumerate}

% ── References ────────────────────────────────────────────────────────
\section*{References}
\begin{itemize}
    \item Potsdam Dataset: \url{https://www.isprs.org/resources/datasets/benchmarks/UrbanSemLab/2d-sem-label-potsdam.aspx}
    \item TensorFlow Guide: \url{https://www.tensorflow.org/tutorials/images/cnn}
    \item Google Colab: \url{https://colab.research.google.com/}
\end{itemize}

\end{document}
